
\lussec 2. QCD 中的梯度流

本文所考虑的理论是规范群为 $\SUn$、含两味或更多有质量夸克
（处于规范群基础表示的多重态）的 QCD。
不过许多结果具有相当的普遍性，并不局限于 QCD。
理论在欧几里得空间中建立，
并照例通过泛函积分量子化。
在本节及下一节中采用维数正规化。
记号约定总结于附录 A。


\lussubsec 2.1 流方程

设 $A_{\mu}(x)$ 为 $\SUn$ 规范势，
$\psi(x)$ 为夸克场，$\psibar(x)$ 为在泛函积分中被积分掉的
反夸克场。后者带有狄拉克、颜色与味道指标，
为简单起见通常将其省略。

杨--米尔斯梯度流把规范场作为一个参数 $t\geq0$ 的函数来演化，
该参数被称为流时间。
从基本规范场出发，
\equation{
  \left.B_{\mu}\right|_{t=0}=A_{\mu},
  \enum
}
随时间演化的场 $B_{\mu}(t,x)$ 由微分方程
\equation{
  \partial_tB_{\mu}=D_{\nu}G_{\nu\mu},
  \enum
  \nexteq{2ex}
  G_{\mu\nu}=\partial_{\mu}B_{\nu}-\partial_{\nu}B_{\mu}+[B_{\mu},B_{\nu}],
  \qquad
  D_{\mu}=\partial_{\mu}+[B_{\mu},\,\cdot\;]
  \enum
}
确定（关于该专题的导引见文献~[\cite{WilsonFlow}]）。

正如第 1 节已经提到的，本文所考虑的流向夸克场的
推广是一个极小化的推广：
规范场的演化保持不变。
然而规范场出现在如下流方程之中\kern1pt\sfn{$\dagger$}{%
任何微分算符 $\Delta$ 的左作用
（在格点场论中则是差分算符的左作用）
通过要求关系式
$\eta^{\dagger}\lvec{\Delta}=(\Delta\eta)^{\dagger}$
对所有复值场 $\eta$ 成立来定义。}
\equation{
  \partial_t\chi=\Delta\chi,
  \qquad
  \partial_t\chibar=\chibar\lvec{\Delta},
  \enum
  \nexteq{2.5ex}
  \Delta=D_{\mu}D_{\mu},\qquad D_{\mu}=\partial_{\mu}+B_{\mu},
  \enum
}
它与初始条件
\equation{
  \left.\chi\right|_{t=0}=\psi,
  \qquad
  \left.\chibar\right|_{t=0}=\psibar,
  \enum
}
一起定义了随时间演化的夸克场与反夸克场
$\chi(t,x)$ 与 $\chibar(t,x)$。
在流方程 (2.4) 中，
规范协变的拉普拉斯算子 $\Delta$
例如可以换成狄拉克算子 $\slash{D}$ 的平方，
或乘以一个比例因子，
但这类修改似乎并无任何好处。

本节引入的夸克场流
与很多年前文献~[\cite{Guesken},\cite{AlexandrouEtAl}]
提出的源光滑化操作类似。
然而相对于这些流行的方法，
存在若干重要差别，
其一是该流在四维时空中运作，
而不是作用在等时超平面上的场上。
此外，流时间是连续变化的，
且规范场并不被固定为基本规范场，
而是与夸克场一同演化。


\lussubsec 2.2 局域复合场

由于流方程是规范协变的，
随时间演化的场在规范变换下与基本场按同样方式变换。
既在时空局域又在流时间局域的规范不变复合场的例子
是如下密度：
\equation{
  E_t(x)=-\frac{1}{2}\tr\{G_{\mu\nu}(t,x)G_{\mu\nu}(t,x)\},
  \enum
  \nexteq{2.5ex}
  S^{rs}_t(x)=\chibar_r(t,x)\chi_s(t,x),
  \qquad
  P^{rs}_t(x)=\chibar_r(t,x)\dirac{5}\chi_s(t,x),
  \enum
}
其中 $r,s$ 为味道指标。

通过初始条件 (2.1)、(2.6)，
随时间演化的场依赖于基本场。
因此从 QCD 泛函积分的观点看，
像密度 (2.7)、(2.8) 这样的复合场是可观测量，
类似于 Wilson 圈或普通的赝标密度等。
下文中，我们关心的是这些场的关联函数，
即由任意流时间的基本场构成的局域场乘积的期望值。


\lussubsec 2.3 微扰论

如文献~[\cite{WilsonFlow},\cite{RenFlow}] 所解释的，
这类关联函数可以直接按规范耦合的幂次展开。
首先将流方程按基本场的幂次求解。
把这些展开代入局域场，
便得到基本场关联函数的线性组合，
后者可用标准 QCD 费曼规则计算。

在微扰论中，
为了避免一些技术上的微妙之处，
将流方程 (2.2) 与 (2.4) 换成
\equation{
  \partial_tB_{\mu}=D_{\nu}G_{\nu\mu}+\alpha_0D_{\mu}\partial_{\nu}B_{\nu},
  \enum
  \nexteq{2.5ex}
  \partial_t\chi=\Delta\chi-
  \alpha_0\partial_{\nu}B_{\nu}\chi,
  \qquad
  \partial_t\chibar=\chibar\lvec{\Delta}+
  \alpha_0\chibar\partial_{\nu}B_{\nu},
  \enum
}
与参数 $\alpha_0>0$ 成比例的各项起到阻尼规范模式的作用，
并且可以通过一个依赖时间的规范变换将其去掉
[\cite{WilsonFlow}]。
因此规范不变场的关联函数不受这些额外项的影响。

将夸克流方程按其积分形式迭代，
便得到随时间演化的夸克场 $\chi(t,x)$
按基本场幂次的展开：
\equation{
  \chi(t,x)=\int\rmd^{D}\kern-1pt y\left\{K_t(x-y)\psi(y)
  +\int_0^t\rmd s\,K_{t-s}(x-y)\Delta'\chi(s,y)\right\},
  \enum
  \nexteq{3.0ex}
  \Delta'=(1-\alpha_0)\partial_{\nu}B_{\nu}
  +2B_{\nu}\partial_{\nu}+B_{\nu}B_{\nu},
  \enum
}
随时间演化的规范场也有相应的积分方程
[\cite{WilsonFlow}]，其中
\equation{
  K_t(z)={\rme^{-z^2/4t}\over(4\pi t)^{D/2}}
  \enum
}
表示 $D$ 维拉普拉斯算子的热核。
在规范耦合的最低阶，
可略去式 (2.11) 中的相互作用项，
于是得到随时间演化夸克场两点函数的表达式
\equation{
  \langle\chi(t,x)\chibar(s,y)\rangle=
  \int{\rmd^{D}\kern-1pt p\over(2\pi)^D}\,
  \rme^{ip(x-y)}
  {\rme^{-(t+s)p^2}\over M_0+i\slash{p}}+\rmO(g_0^2)
  \enum
}
其中 $g_0$ 与 $M_0$ 分别为裸耦合与（对角的）夸克质量矩阵。
由这些方程可以明显看出夸克流的光滑化特征。
此外可见，
领头阶下的光滑化范围 $\sqrt{8t}$ 与规范场情形相同。

在单圈微扰论中，
有 $8$ 个自能图对两点函数 (2.14) 有贡献。
这些图中的顶点来自 QCD 作用量
以及积分方程 (2.11) 的迭代。
按照文献~[\cite{RenFlow}] 的思路，
容易算出各图的紫外发散部分，
并且随即发现：只需照常重正化夸克质量，
并按下式重正化场，
两点函数即可被重正化：
\equation{
  \chi=Z_{\chi}^{-1/2}\ren{\chi},
  \qquad
  \chibar=\ren{\chibar}Z_{\chi}^{-1/2}.
  \enum
}
在 $D=4-2\eps$ 维的 $\MSbar$ 方案中，
对该场重正化常数的计算给出
\equation{
  Z_{\chi}=1+{3\CF\over16\pi^2\eps}g^2+\rmO(g^4),
  \qquad
  \CF={N^2-1\over2N},
  \enum
}
其中 $g$ 为重正化耦合。


\lussubsec 2.4 非零流时间下的夸克凝聚

在非零流时间下，
积分 (2.14) 中的大动量被指数压低，
因而当 $(t,x)\to(s,y)$ 时夸克两点函数没有奇异性。
短距离奇异性的缺失是正流时间下关联函数的一般特征，
它源于流方程的光滑化性质。
特别地，``随时间演化的夸克凝聚''
\equation{
  \Sigt^{rr}=-\langle S^{rr}_t(x)\rangle
  \enum
}
不需要加性重正化。

在微扰论中，
重正化的凝聚
\equation{
  \rens{\Sigma}{t}^{rr}=Z_{\chi}\Sigt^{rr}
  \enum
}
可以按重正化耦合 $g$ 的幂次算出，
其系数依赖于重正化夸克质量 $\mR{r},\mR{s},\ldots$
与流时间 $t$。
四维下的领头阶项由下式给出：
\equation{
  \left.\rens{\Sigma}{t}^{rr}\right|_{g=0}=
  {2N\mR{r}\over(4\pi)^2t}
  \int_0^{\infty}\rmd v\,{\rme^{-vz}\over(1+v)^2},
  \qquad z=2t\mR{r}^2,
  \enum
}
因而在手征极限下为零，
这与手征对称性只能在非微扰层面发生自发破缺这一事实相一致。
